% ============================================================
% CHAPITRE 4
% ============================================================
\chapter{Sprint 3 : Commandes et paiement}

\section{Introduction}

Transformer le visiteur affamé en acheteur confirmé, tel est l'enjeu du Sprint 3. Nous y décortiquons comment l'architecture numérique digère un panier virtuel pour le changer en ordre de préparation en cuisine. La particularité de cette phase est son arrimage au marché tunisien : nous avons sciemment privilégié le paiement de la main à la main (cash-on-delivery) pour correspondre aux habitudes de la clientèle d'Assiette Gala, ce qui conditionne tout le code lié à la facturation.

% --------------------------------------------------
\section{Backlog du sprint}

\begin{table}[H]
    \centering
    \caption{Backlog du sprint 3}
    \label{tab:backlog_sprint3}
    \renewcommand{\arraystretch}{1.3}
    \begin{tabular}{|c|p{8cm}|c|}
        \hline
        \textbf{ID} & \textbf{User Story} & \textbf{Points} \\
        \hline
        US-14 & En tant que client, je souhaite ajouter des plats à mon panier. & 3 \\
        \hline
        US-15 & En tant que client, je souhaite passer une commande. & 8 \\
        \hline
        US-16 & En tant que client, je souhaite passer commande avec paiement en espèces. & 5 \\
        \hline
        US-17 & En tant que client, je souhaite suivre l'état de ma commande. & 3 \\
        \hline
        US-18 & En tant que client, je souhaite télécharger la facture de ma commande. & 5 \\
        \hline
        US-19 & En tant qu'administrateur, je souhaite gérer les commandes. & 5 \\
        \hline
    \end{tabular}
\end{table}

% --------------------------------------------------
\section{Conception}

\subsection{Diagramme de cas d'utilisation}

\begin{figure}[H]
    \centering
    \includegraphics[width=0.85\textwidth]{uc_sprint3.png}
    \caption{Diagramme de cas d'utilisation du sprint 3}
    \label{fig:uc_sprint3}
\end{figure}

\subsection{Description textuelle des cas d'utilisation}

\subsubsection{Cas d'utilisation : Passer une commande}

\begin{table}[H]
    \centering
    \caption{Description textuelle --- Passer une commande}
    \label{tab:desc_passer_commande}
    \renewcommand{\arraystretch}{1.3}
    \begin{tabular}{|p{4cm}|p{8.5cm}|}
        \hline
        \textbf{Cas d'utilisation} & Passer une commande \\
        \hline
        \textbf{Acteur principal} & Client \\
        \hline
        \textbf{Précondition} & Le client est authentifié et son panier contient au moins un plat \\
        \hline
        \textbf{Postcondition} & La commande est créée avec le statut PENDING \\
        \hline
        \textbf{Scénario principal} &
        \begin{enumerate}
            \item L'appétit aiguisé, l'acheteur clique sur l'icône de sa besace (panier)
            \item Présentation d'un bref ticket reprenant les repas et la TVA
            \item Renseignement de l'adresse de livraison (domicile ou travail)
            \item Indication de la volonté d'être livré ou de récupérer sur place
            \item Le bouton de validation déclenche l'envoi vers le backend
            \item PostgreSQL réceptionne et fige le dossier en statut "En attente"
            \item Affichage d'un reçu d'écran avec l'exigence du montant en espèce
        \end{enumerate} \\
        \hline
        \textbf{Scénario alternatif} &
        \begin{itemize}
            \item[3a.] Client sans adresse connue : forçage de la saisie GPS/texte
            \item[5a.] Épuisement d'un plat dans la seconde : rejet avec message désolé
        \end{itemize} \\
        \hline
    \end{tabular}
\end{table}

\subsubsection{Cas d'utilisation : Confirmer le paiement en espèces}

\begin{table}[H]
    \centering
    \caption{Description textuelle --- Paiement en espèces}
    \label{tab:desc_paiement}
    \renewcommand{\arraystretch}{1.3}
    \begin{tabular}{|p{4cm}|p{8.5cm}|}
        \hline
        \textbf{Cas d'utilisation} & Confirmer le paiement en espèces \\
        \hline
        \textbf{Acteur principal} & Administrateur \\
        \hline
        \textbf{Précondition} & Une commande a été créée avec le statut PENDING \\
        \hline
        \textbf{Postcondition} & Le paiement est confirmé, le statut devient COMPLETED \\
        \hline
        \textbf{Scénario principal} &
        \begin{enumerate}
            \item Le livreur d'Assiette Gala perçoit la liasse de billets chez le client
            \item Le gérant ouvre la tour de contrôle de son interface
            \item Repérage de la transaction incriminée via son numéro UUID
            \item Bascule manuelle du marqueur de "En attente" vers "Payé/Livré"
            \item Le système verrouille définitivement le cycle comptable
        \end{enumerate} \\
        \hline
        \textbf{Scénario alternatif} &
        \begin{itemize}
            \item[1a.] Rétractation du client à la porte : marquage "Annulé" et perte de nourriture
        \end{itemize} \\
        \hline
    \end{tabular}
\end{table}

\subsubsection{Cas d'utilisation : Télécharger la facture}

\begin{table}[H]
    \centering
    \caption{Description textuelle --- Télécharger la facture}
    \label{tab:desc_facture}
    \renewcommand{\arraystretch}{1.3}
    \begin{tabular}{|p{4cm}|p{8.5cm}|}
        \hline
        \textbf{Cas d'utilisation} & Télécharger la facture \\
        \hline
        \textbf{Acteur principal} & Client \\
        \hline
        \textbf{Précondition} & La commande a le statut COMPLETED \\
        \hline
        \textbf{Postcondition} & Le fichier PDF de la facture est téléchargé \\
        \hline
        \textbf{Scénario principal} &
        \begin{enumerate}
            \item Accès au tableau de bord privé du consommateur
            \item Clic sur le pictogramme d'impression du ticket
            \item Le module PDFKit dessine dynamiquement la grille comptable
            \item Le transfert binaire est offert en téléchargement direct (PDF)
        \end{enumerate} \\
        \hline
    \end{tabular}
\end{table}

\subsection{Diagrammes de séquence}

\subsubsection{Séquence : Processus de commande et paiement}

\begin{figure}[!htbp]
    \centering
    \includegraphics[width=0.9\textwidth]{seq_commande_paiement.png}
    \caption{Diagramme de séquence --- Commande et paiement}
    \label{fig:seq_commande_paiement}
\end{figure}

\subsubsection{Séquence : Génération de facture PDF}

\begin{figure}[!htbp]
    \centering
    \includegraphics[width=0.9\textwidth]{seq_facture.png}
    \caption{Diagramme de séquence --- Génération de facture}
    \label{fig:seq_facture}
\end{figure}

\subsection{Diagramme de classes}

\begin{figure}[!htbp]
    \centering
    \includegraphics[width=0.92\textwidth,height=0.72\textheight,keepaspectratio]{class_sprint3.png}
    \caption{Diagramme de classes du service de commandes}
    \label{fig:class_sprint3}
\end{figure}

Le modèle de données du service de commandes repose sur trois entités principales :

\textbf{Le modèle \texttt{Order}} contient :
\begin{itemize}
    \item \texttt{id} : identifiant unique (UUID)
    \item \texttt{userId} : identifiant du client ayant passé la commande
    \item \texttt{status} : état de la commande (\texttt{PENDING}, \texttt{PROCESSING}, \texttt{COMPLETED}, \texttt{CANCELLED}, \texttt{FAILED})
    \item \texttt{totalAmount} : montant total de la commande (type décimal)
    \item \texttt{deliveryMode} : mode de livraison (\texttt{DELIVERY}, \texttt{PICKUP})
    \item \texttt{deliveryAddress} : adresse de livraison (JSON)
    \item \texttt{items} : collection de lignes de commande
\end{itemize}

\textbf{Le modèle \texttt{OrderItem}} contient :
\begin{itemize}
    \item \texttt{id} : identifiant unique
    \item \texttt{platId} : identifiant du plat commandé
    \item \texttt{platName} : nom du plat (dénormalisé pour l'historique)
    \item \texttt{platPrice} : prix unitaire au moment de la commande
    \item \texttt{quantity} : quantité commandée
    \item \texttt{orderId} : clé étrangère vers \texttt{Order}
\end{itemize}

\subsection{Diagramme d'états-transitions}

Le diagramme d'états-transitions (figure~\ref{fig:state_sprint3}) illustre le cycle de vie complet d'une commande, depuis sa création jusqu'à sa livraison ou son annulation.

\begin{figure}[!htbp]
    \centering
    \includegraphics[width=0.9\textwidth,height=0.72\textheight,keepaspectratio]{state_sprint3_order.png}
    \caption{Diagramme d'états-transitions --- Cycle de vie d'une commande}
    \label{fig:state_sprint3}
\end{figure}

% --------------------------------------------------
\section{Réalisation}

\subsection{Points d'entrée de l'API}

\begin{table}[H]
    \centering
    \caption{Endpoints REST du service de commandes}
    \label{tab:endpoints_order}
    \renewcommand{\arraystretch}{1.2}
    \begin{tabular}{|l|p{5cm}|p{4.5cm}|}
        \hline
        \textbf{Méthode} & \textbf{Route} & \textbf{Description} \\
        \hline
        POST & /api/orders & Création d'une commande \\
        \hline
        GET & /api/orders & Commandes de l'utilisateur \\
        \hline
        GET & /api/orders/:id & Détails d'une commande \\
        \hline
        GET & /api/orders/:id/invoice & Téléchargement facture PDF \\
        \hline
        GET & /api/orders/admin/all & Toutes les commandes (administrateur) \\
        \hline
        PUT & /api/orders/admin/:id/status & Mise à jour du statut (administrateur) \\
        \hline
    \end{tabular}
\end{table}

% Schéma de la base de données retiré : le diagramme de classes (figure~\ref{fig:class_sprint3}) suffit à modéliser la structure des données.



\subsection{Gestion du paiement en espèces}

Dans le tissu commercial sfaxien, la grande majorité des transactions traiteur s'effectuent en liquide au moment de la réception du plateau. Nous avons donc bâti un circuit qui épouse cette habitude locale plutôt que de la contredire :

\begin{enumerate}
    \item \textbf{Dépôt de la demande :} l'acheteur finalise son panier en ligne ; la commande atterrit dans PostgreSQL avec l'étiquette \texttt{PENDING}.
    \item \textbf{Acquittement par le gérant :} M. Daoud ouvre le panneau d'administration, vérifie la faisabilité (stock, disponibilité de la date) et bascule le dossier en \texttt{PROCESSING}.
    \item \textbf{Règlement physique :} le livreur perçoit la liasse de billets au seuil du client.
    \item \textbf{Clôture comptable :} un clic fait passer le statut à \texttt{COMPLETED}, verrouillant la traçabilité financière.
\end{enumerate}

Ce modèle supprime toute friction liée aux passerelles bancaires et élimine les commissions tierces, tout en conservant un audit trail complet dans la base de données.

\subsection{Génération de factures PDF}

Chaque transaction validée peut donner lieu à une pièce comptable imprimable, forgée à la volée par le module \texttt{PDFKit}. Le document embarque :

\begin{itemize}
    \item Le cartouche Assiette Gala (logo, adresse Route Mahdia km 1, téléphone)
    \item Les coordonnées de l'acheteur
    \item La référence UUID et la date de captation
    \item La ventilation détaillée (dénomination, quantité, prix HT, sous-total)
    \item Le total TTC incluant la TVA à 19\,\%
    \item La mention de l'état du règlement
\end{itemize}

Le fichier binaire est assemblé en mémoire vive puis projeté vers le navigateur avec l'entête \texttt{Content-Type: application/pdf}, déclenchant le téléchargement immédiat.

\subsection{Interfaces réalisées}

\subsubsection{Panier d'achat}

\begin{figure}[H]
    \centering
    \includegraphics[width=0.85\textwidth]{ui_panier.png}
    \caption{Interface du panier d'achat}
    \label{fig:ui_panier}
\end{figure}

Le panier récapitule la sélection du gourmand : chaque ligne dévoile la vignette du mets, son tarif unitaire et un compteur incrémentable. Le bas de page totalise la facture prévisionnelle, invitant l'acheteur à modifier ou élaguer avant de poursuivre.

\subsubsection{Page de validation de commande}

\begin{figure}[H]
    \centering
    \includegraphics[width=0.85\textwidth]{ui_checkout.png}
    \caption{Interface de validation de commande}
    \label{fig:ui_checkout}
\end{figure}

Cet écran rassemble en un coup d'\oe{}il le résumé du panier, le choix de l'adresse de livraison et le mode de réception (livraison à domicile ou retrait Route Mahdia). Le total TTC en dinars tunisiens, calculé avec une TVA de 19\,\%, est mis en évidence.

\subsubsection{Historique des commandes}

\begin{figure}[H]
    \centering
    \includegraphics[width=0.85\textwidth]{ui_historique_commandes.png}
    \caption{Interface de l'historique des commandes}
    \label{fig:ui_historique}
\end{figure}

Le journal personnel du client empile ses achats passés du plus récent au plus ancien : date, montant déboursé et pastille colorée (jaune = en attente, vert = livré, rouge = annulé). Un lien ouvre la fiche détaillée de chaque dossier.

\subsubsection{Gestion des commandes (Administrateur)}

\begin{figure}[H]
    \centering
    \includegraphics[width=0.9\textwidth]{ui_admin_commandes.png}
    \caption{Interface d'administration des commandes}
    \label{fig:ui_admin_commandes}
\end{figure}

Le poste de contrôle de M. Daoud dresse un inventaire filtré par état. D'un clic, il fait progresser une commande de \texttt{PROCESSING} vers \texttt{COMPLETED}, et peut ouvrir le détail complet (coordonnées, plats, montant) de chaque opération.

% --------------------------------------------------
\section{Tests fonctionnels}

\begin{table}[H]
    \centering
    \caption{Tests fonctionnels --- Sprint 3}
    \label{tab:tests_sprint3}
    \renewcommand{\arraystretch}{1.3}
    \begin{tabular}{|c|p{3cm}|p{3.5cm}|p{2.5cm}|}
        \hline
        \textbf{N°} & \textbf{Cas de test} & \textbf{Comportement attendu} & \textbf{Résultat} \\
        \hline
        1 & Ajout d'un plat au panier & Plat ajouté, quantité et total mis à jour & Réussi \\
        \hline
        2 & Modification de la quantité dans le panier & Sous-total et total recalculés & Réussi \\
        \hline
        3 & Suppression d'un article du panier & Article retiré, total recalculé & Réussi \\
        \hline
        4 & Validation de la commande & Commande créée en base, statut PENDING & Réussi \\
        \hline
        5 & Paiement en espèces confirmé & Statut mis à jour à COMPLETED par l'administrateur & Réussi \\
        \hline
        6 & Annulation de commande & Statut mis à jour à CANCELLED, message affiché & Réussi \\
        \hline
        7 & Téléchargement de facture PDF & PDF généré contenant tous les détails & Réussi \\
        \hline
        8 & Consultation de l'historique & Liste des commandes avec statuts corrects & Réussi \\
        \hline
        9 & Mise à jour du statut par l'administrateur & Statut modifié en base et sur l'interface & Réussi \\
        \hline
        10 & Commande sans authentification & Redirection vers la page de connexion & Réussi \\
        \hline
    \end{tabular}
\end{table}

% --------------------------------------------------
\section{Conclusion}

À l'issue de cet acte de programmation, Assiette Gala jouit d'une caisse enregistreuse immatérielle connectée directement à sa cuisine. Le traitement du panier, calqué sur le comportement tunisien du "paiement à l'arrivée", garantit une conversion commerciale maximale car débarrassé de la friction des cartes bleues. Fort de cette ossature financière, le Sprint 4 peut s'atteler à un aspect encore plus majestueux : l'organisation complète des noces et cérémonies.
